\section{Introduction}\label{sec:introduction}

Modern homes more and more depend on wireless devices such as smart speakers,
video doorbells, and thermostats that connect over Bluetooth, Wi-Fi, or a
cellular network. The software controlling these radios is typically large,
complex, and closed-source, which makes vulnerabilities inevitable. Because these
systems receive signals over the air, anyone nearby can attempt to exploit a
weakness without a password or prior access. Proximity alone is enough. The risk
is not small. Vast numbers of these radios are deployed inside trusted home
environments and are rarely updated, so a single flaw can become a persistent,
unauthorised way in.

The standard way to surface such flaws is \emph{fuzzing}, which sends a target
many malformed or out-of-order packets and watches for anything that crashes it.
Fuzzing is effective. The \textsc{SweynTooth} campaign discovered twelve new
Bluetooth Low Energy vulnerabilities across seven vendors, while
\textsc{BrakTooth} found another sixteen in Bluetooth Classic
devices~\cite{sweyntooth2020,braktooth2022}. Simply detecting a crash, though,
does not give a developer enough to fix the issue.

What a developer needs is a \emph{Proof-of-Concept} (PoC): the minimal set of
packets that reliably triggers the fault and pins down the specific message and
code path involved. Distilling that sequence is the slow part. A campaign may
send tens of thousands of packets, yet when a crash occurs it typically records
only a terse memory dump, with no indication of which packets mattered.
Diagnosing each crash by hand stalls a triage team.
\textsc{AirBugCatcher}~\cite{airbugcatcher2024} is currently the leading tool for
automating this last step. Seeing where it is forced to compromise is the way
into our proposal.

The main challenge comes from the unpredictable nature of wireless communication.
Re-sending the same packets does not always yield the same result, because
devices keep hidden internal state and the channel adds noise. Every reproduction
tool depends on an \emph{oracle}~\cite{oracleproblem2015}, the component that
decides whether a sequence triggered the bug, yet that oracle can return
inconsistent results, which makes accurate sequence minimisation hard. Delta
debugging~\cite{ddmin2002} is a proven technique for isolating failure-inducing
inputs, but it assumes the oracle is always consistent. When that assumption
fails, the method can discard critical packets and converge on an incorrect or
empty result. \textsc{AirBugCatcher} copes by restricting recipes to a small,
fixed size and testing combinations within a time budget. It averages over the
noise rather than eliminating it, and so gives up the efficiency and minimality
that delta debugging is meant to provide.

We introduce \emph{Robust Delta-Debugging} (RDD), which restores the
stable-verdict assumption in software and so lets delta debugging work even on
noisy wireless targets. Rather than circumventing the problem, RDD repairs the
core issue the algorithm depends on. It uses a sequential statistical test, a
truncated Sequential Probability Ratio Test (SPRT), to pose the noisy
``did it crash?'' question repeatedly and aggregate the answers into a single
reliable verdict with a controlled false-positive rate. That verdict is only as
reliable as the tool's ability to tell whether two crash dumps come from the same
bug, so RDD adds a two-tier identity oracle: a fast, rule-based matcher for most
cases, and a single open-weight model for the ambiguous ones. A robust minimiser
then drives the search even where a packet conceals the crash. Nothing is
reported until an independent check re-confirms the recipe by reproducing the
crash, which keeps a wrongly credited PoC improbable by the same tunable margin.

This report makes four main contributions. The first identifies the oracle,
rather than the search strategy, as the binding constraint in wireless
reproduction. The second introduces RDD, a sound pipeline that reinstates the
stable-verdict assumption through a statistically stabilised oracle. The third
evaluates the approach with a prototype, measured against a faithful
re-implementation of \textsc{AirBugCatcher} on real Zephyr controller bugs and
under heavy synthetic stress, recovering 0.863 of reproductions at zero false
credit where the baseline manages 0.759 at 0.137, and 0.916 versus 0.282
under tougher conditions. % B1 headtohead.json + B2 resilience.json OVERALL
The fourth turns the method's structural limits into a focused research agenda.
The work is academically oriented. Its results are measured on public, open-source
targets, so anyone can reproduce them, and the deliverable is a supported research
direction with clear next steps, not a finished product. Section~\ref{sec:background} gives the background.
Section~\ref{sec:our_approach} presents RDD. Section~\ref{sec:evaluation} provides
the evidence and proposed experiments. Section~\ref{sec:related_work} reviews the
related work. Section~\ref{sec:conclusions} concludes.
